\paragraph{The metric comparison, with its exact additive budget}
Put $t=32s$, $(a_1,a_2,a_3,a_4)=(0,32,76832,-343)$ and
$P(t)=\prod_i(t-a_i)$. The sections
\[
 q_0=(dt)^2/P(t),\qquad q_1=tq_0
\]
are holomorphic bicanonical sections on $X$, with no common zero. At a
finite branch point use $t-a_i=z^2$; at infinity use $t=z^{-2}$.
At infinity $q_0$ vanishes to order two and $q_1$ does not vanish.

For each prime let $\delta_p=\max_{i<j}v_p(a_i-a_j)$, so
\[
 J=\prod_p p^{\delta_p}=2^{10}3^2 5^3 7^4=2765952000.
\]
Both $J^2q_i$ extend across the maximal regular stack. Indeed, at the
generic point of a vertical component semistability makes $p$ a
uniformizer, and the relative differential satisfies $v(du)\ge v(u)$.
Writing $r_i=v(t-a_i)$ gives $v(dt)\ge\max r_i$. If $v(t)<0$, all $r_i$
equal $v(t)$ and both sections have nonnegative valuation. Otherwise all
$r_i\ge0$ and at most one exceeds $\delta_p$. In either case
$2\max r_i-\sum r_i\ge-2\delta_p$. This proves the required vertical
integrality; regularity handles codimension two.

For the archimedean estimate write $r_i=|t-a_i|$, $U_0=1+\sum_i r_i$ and
\[
 ds_\rho^2=\frac{U_0}{\prod_i r_i}|dt|^2,\qquad
 C_{\rm met}=16(12\cdot76832+9)^3=12540153541944538512.
\]
Its pullback is smooth and positive at every branch point and infinity.
The curvature formula, $\Delta U_0=\sum_i r_i^{-1}$ and
$|\nabla U_0|^2\le16$ give
\[
 -K_\rho\ge\frac{\sum_i\prod_{j\ne i}r_j}{2U_0^3}.
\]
When $|t|\le2\cdot76832+2$, at most one $r_i<1/2$, and
$U_0\le12\cdot76832+9$. The displayed bound is at least
$C_{\rm met}^{-1}$. Outside this disk, $|t|/2\le r_i\le3|t|/2$
and $U_0\le7|t|$, giving the stronger bound $1/1372$.
The maximum principle against the curvature $-1$ metric implies
$ds_\rho^2/C_{\rm met}\le ds_{\rm hyp}^2$. Yuan's Petersson norm is
$\|d\tau\|=2\operatorname{Im}\tau$. Consequently
\[
 \|q_0\|_{\rm Pet}\le4C_{\rm met}/U_0,\qquad
 \|q_1\|_{\rm Pet}\le4C_{\rm met}|t|/U_0\le4C_{\rm met}.
\]

Yuan's isometry \cite[Theorem 1.1 and metric definitions, pp.~2--3]{Yuan2024} identifies the square of the Hodge bundle with
$210\omega_{\mathscr X}^{\otimes2}$ with the induced metric. The Hodge
height in his normalization equals the original stable Faltings height
plus $\log\pi$ in dimension two. Thus
\[
 h_{\bar\omega^{\otimes2}}(y)=2h_F(A_y)+2\log\pi+\log210.
\]
The sign of $\log210$ is positive: the sublattice $210\omega^{\otimes2}$
has height $h_{\bar\omega^{\otimes2}}-\log210$. Evaluate on an extension
where the point extends to the proper maximal stack; its universal
abelian scheme computes stable height. The integral sections and the
supremum estimate yield
\[
 h(s)\le2h_F(A_y)+2\log\pi+\log210+2\log J
                       +\log(4C_{\rm met})+\log32.
\]
The integer comparison
$210J^2(4C_{\rm met})32\cdot16<2^{160}$, together with $\pi<4$ and
$\log2<3/4$, bounds the additive term by $120<128$. Therefore
\[
 h(s)\le2\max(0,h_F(A_y))+128.
\tag{G}\label{eq:h5-G}
\]

